\documentclass[10pt,a4paper]{amsart}
\usepackage[margin=19mm]{geometry}
\usepackage{amsmath,amssymb,mathtools}
\usepackage[scr=rsfs,cal=euler]{mathalfa}
\usepackage{xfrac}
\usepackage[unicode,hidelinks,bookmarks=false]{hyperref}
\newcommand{\ie}{i.e.\ }
\newcommand{\cf}{cf.\ }
\newcommand{\resp}{respectively\ }
\newcommand{\Subsection}{\subsection}
\providecommand{\nobreakdash}{\nobreak\hbox{-}}
\setlength{\emergencystretch}{2em}
\multlinegap0pt
\usepackage{amssymb}
\usepackage{braket}
\usepackage{tikz}
\newtheorem{definition}{Definition}
\newtheorem{lemma}{Lemma}
\newcommand{\MC}{\mathrm{MC}}
\title{Bigraded path homology and the magnitude-path spectral sequence}
\author{Richard Hepworth, Emily Roff}
\date{Source: arXiv:2404.06689v3 (CC BY 4.0)}
\begin{document}
\maketitle

\noindent\emph{Source and licence.} Bigraded path homology and the magnitude-path spectral sequence. Version arXiv:2404.06689v3 (CC BY 4.0), distributed by arXiv under the Creative Commons Attribution 4.0 International licence, http://creativecommons.org/licenses/by/4.0/, as recorded in the arXiv licence metadata for this exact version and shown on the abstract page https://arxiv.org/abs/2404.06689v3. Full licence notice: \texttt{licenses/arxiv-2404.06689-CC-BY-4.0.txt}.\par\smallskip

\noindent\textit{Editorial scope.} Counted unit: the article's lemma A-acyclic (source0.tex lines 2197--2202). The article's complete original proof of it is delivered with it (source0.tex lines 2204--2399, one \texttt{proof} environment of 196 lines). No mathematics is written, repaired or re-derived: the statement and the proof are the article's own lines, in the article's own notation, and no retained line is substituted or re-flowed.  Retained uncounted as the source-local context the proof needs: lines 2154--2162.  Nothing was omitted from any retained span, so the package carries no omission record.  No retained line contains a citation, so the wrapper carries no bibliography and no bibliography entry.  No retained line contains a figure environment, an image-inclusion command or drawing code, so no image is delivered and the package declares no asset dependency.  Editorial formatting only, in addition: an AMS \texttt{amsart} preamble that restates the article's own declarations and macros used by the retained lines, namely the environments \texttt{definition}, \texttt{lemma} on separate counters, together with the standard packages those lines need.  The retained environments print in this document as Definition 1, Lemma 1 in order of appearance, whereas the article prints the same items with the numbering of its own full document.  The complete native source of the article as distributed in the arXiv e-print for this exact version is bundled unchanged as \texttt{sources/157471/source0.tex}.
\begin{definition}\label{definition-A}
    Let $A\hookrightarrow X$ be a cofibration,
    let $\ell\geq 0$, let $x\in X\setminus A$
    and let $a\in A$.
    Define $A_{\ast,\ell}(a,x)$ to be the subcomplex of
    $\MC_{\ast,\ell}(X)$ spanned by those tuples
    $(x_0,\ldots,x_k)$ 
    for which $x_0=a$, $x_k=x$ and $x_1,\ldots,x_{k-1}\in A$.
\end{definition}

\begin{lemma}\label{A-acyclic}
    Suppose we are in the situation of 
    Definition~\ref{definition-A},
    and that at least one of $a=\pi(x)$ and $\ell=d(a,x)$ fails.
    Then $A_{\ast,\ell}(a,x)$ is acyclic.
\end{lemma}

\begin{proof}
    Define a map
    \[
    	s\colon A_{\ast,\ell}(a,x)
            \longrightarrow A_{\ast+1,l}(a,x)
    \]
    by
    \[
    	s(x_0,\ldots,x_k) 
    	=
    	\begin{cases}
    		(-1)^k(x_0,\ldots,x_{k-1},\pi(x_k),x_k)
    		& \mathrm{if}\ x_{k-1}\neq\pi(x_k),
    		\\
    		0
    		&
    		\mathrm{if}\ x_{k-1}=\pi(x_k).
    	\end{cases}
    \]
    We claim that $s$ is a chain homotopy from the identity
    of $A_{\ast,\ell}(a,x)$ to the zero map,
    or in other words that
    \begin{equation}\label{chain-homotopy}
        \partial\circ s + s\circ\partial=\mathrm{Id}.
    \end{equation}
    This immediately shows that the homology 
    of $A_{\ast,\ell}(a,x)$ is trivial, and the result
    follows.
    
    Equation~\eqref{chain-homotopy} is equivalent to the claim
    that 
    \begin{equation}\label{chain-homotopy-expanded}
    	\sum_{i=1}^k (-1)^i \partial_i s(x_0,\ldots,x_k)
    	+ \sum_{i=1}^{k-1}(-1)^i s\partial_i(x_0,\ldots,x_k)
    	=(x_0,\ldots,x_k)
    \end{equation}
    for each generator $(x_0,\ldots,x_k)$ of $A_{\ast,\ell}(a,x)$.

    To begin we consider the case $k=1$.
    In this case  the only possible generator is $(a,x)$, 
    but this is only present when $\ell = d(a,x)$.
    So by our assumption that we do not have both
    $a=\pi(x)$ and $\ell = d(a,x)$, we must have $a\neq\pi(x)$.
    Then~\eqref{chain-homotopy-expanded} becomes
    $
        -\partial_1 s(a,x)
        =
        (a,x),
    $ 
    which follows from the definition of $s$. 

    For the remainder of the proof we consider the case $k\geq 2$.
    We have $\partial_i s = -s\partial_i$ for $1\leq i\leq k-2$,
    since then $\partial_i$ affects only the first $k-1$ terms
    of a tuple, while $s$ affects, and depends upon, only the 
    remaining terms.
    So~\eqref{chain-homotopy-expanded} reduces to the claim
    that:
    \begin{equation}\label{chain-homotopy-simplified}
        \begin{split}
    	&(-1)^{k-1}\ \partial_{k-1}s(x_0,\ldots,x_k)
            \\
    	+
    	&(-1)^{k\phantom{-1}}\ \partial_{k\phantom{-1}} s(x_0,\ldots,x_k)
            \\
    	+&(-1)^{k-1}\ s\partial_{k-1}(x_0,\ldots,x_k)
            =
            (x_0,\ldots,x_k)
        \end{split}
    \end{equation}
    We now divide into the following cases:
    \begin{itemize}
        \item
        Assume that $x_{k-1}=\pi(x_k)$.
        
        Here we see immediately that  
        the first two terms of~\eqref{chain-homotopy-simplified} 
        vanish by the definition of $s$.
        For the third term, we have
        \begin{align*}
            (-1)^{k-1}s\partial_{k-1}(x_0,\ldots,x_k)
            &=
            (-1)^{k-1}s(x_0,\ldots,x_{k-2},x_k)
            \\
            &=
            (x_0,\ldots,x_{k-2},\pi(x_k),x_k)
            \\
            &=
            (x_0,\ldots,x_k)
        \end{align*}
        as required.
        The first of these equations holds because 
        $d(x_{k-2},x_{k-1})+d(x_{k-1},x_k) = d(x_{k-2},x_k)$
        by the defining property of $\pi(x_k)$,
        and the second holds since  $x_{k-2}\neq x_{k-1}=\pi(x_k)$.
        
        \item
        Assume that $x_{k-1}\neq \pi(x_k)$
        and
        $d(x_{k-2},x_{k-1})+d(x_{k-1},x_k) > d(x_{k-2},x_k)$.
        
        Applying the defining property of $\pi(x_k)$ to the 
        assumed inequality, we find that we also have 
        $
            d(x_{k-2},x_{k-1})+d(x_{k-1},\pi(x_k))
            > d(x_{k-2},\pi(x_k))
        $
        so that we have $\partial_{k-1}(x_0,\ldots,x_{k-1},\pi(x_k),x_k)=0$
        and the first term of~\eqref{chain-homotopy-simplified}
        vanishes.
        The assumed inequality also shows that 
        $\partial_{k-1}(x_0,\ldots,x_k)=0$
        so that the third term 
        of~\eqref{chain-homotopy-simplified} vanishes.
        It remains to show that the second term 
        of~\eqref{chain-homotopy-simplified} is 
        $(x_0,\ldots,x_k)$. 
        Thanks to the assumption $x_{k-1}\neq \pi(x_k)$
        we have
        \[
            (-1)^k \partial_k s(x_0,\ldots,x_k)
            =
            \partial_k (x_0,\ldots,x_{k-1},\pi(x_k),x_k)
            =
            (x_0,\ldots,x_{k-1},x_k)
        \]
        as required.
        
        \item
        Assume that 
        $x_{k-1}\neq \pi(x_k)$
        and 
        $d(x_{k-2},x_{k-1})+d(x_{k-1},x_k) = d(x_{k-2},x_k)$.
        
        By applying the defining property of $\pi(x_k)$ to
        the assumed equation, we get
        $d(x_{k-2},x_{k-1})+d(x_{k-1},\pi(x_k)) 
        = d(x_{k-2},\pi(x_k))$,
        and in particular $x_{k-2}\neq \pi(x_k)$.
        The first term of~\eqref{chain-homotopy-simplified}
        is given by:
        \begin{align*}
    	(-1)^{k-1}\partial_{k-1}s(x_0,\ldots,x_k)
            &=
    	-\partial_{k-1}(x_0,\ldots,x_{k-1},\pi(x_k),x_k)
            \\
            &=
            -(x_0,\ldots,x_{k-2},\pi(x_k),x_k),
        \end{align*}
        where the first equation used the assumption
        $x_{k-1}\neq\pi(x_k)$, 
        and the second equation used the fact that
        $d(x_{k-2},x_{k-1})+d(x_{k-1},\pi(x_k))
        = d(x_{k-2},\pi(x_k))$ and $x_{k-2}\neq\pi(x_k)$.
        The second term of~\eqref{chain-homotopy-simplified} is 
        \begin{align*}
    	(-1)^k \partial_k s(x_0,\ldots,x_k)
            &=
            \partial_k (x_0,\ldots,x_{k-1},\pi(x_k),x_k)
            \\
            &=
    	(x_0,\ldots,x_{k-1},x_k),
        \end{align*}
        where again in the first equation we used the assumption
        $x_{k-1}\neq\pi(x_k)$, and in the second equation
        we used the defining property of 
        $\pi(x_k)$ to see that
        $d(x_{k-1},\pi(x_k))+d(\pi(x_k),x_k) = d(x_{k-1},x_k)$.
        The third term of~\eqref{chain-homotopy-simplified} is  
        \begin{align*}
    	(-1)^{k-1}s\partial_{k-1}(x_0,\ldots,x_k)
            &=
    	(-1)^{k-1}s(x_0,\ldots,x_{k-2},x_k)
            \\
            &=
    	(x_0,\ldots,x_{k-2},\pi(x_k),x_k)
        \end{align*} 
        where in the first equation we used the assumption
        that
        $d(x_{k-2},x_{k-1})+d(x_{k-1},x_k) = d(x_{k-2},x_k)$,
        and in the second we used the fact $x_{k-2}\neq\pi(x_k)$.
        So altogether, the left hand side 
        of~\eqref{chain-homotopy-simplified} is
        \[
            -(x_0,\ldots,x_{k-2},\pi(x_k),x_k)
            +
    	(x_0,\ldots,x_{k-1},x_k),
            +
    	(x_0,\ldots,x_{k-2},\pi(x_k),x_k)
        \]
        which is precisely $(x_0,\ldots,x_k)$ as required.
    \end{itemize}
    In all three cases above, the equation
    \eqref{chain-homotopy-simplified} holds.
    This completes the proof.
\end{proof}

\end{document}
